% GENERATED FILE. Edit canonical lessons, then run npm run generate:books.
% canonical-source-hash: fnv1a64:9935343703e4248e
% canonical-lessons: FR-C528-balayer, FR-C528-epousseter, FR-C528-frotter, FR-C528-eponge, FR-C528-serpilliere

\chapter{Faire le ménage}
\label{ch:fr-faire-le-menage}
\hlchaptermodality{528}

\begin{chapteropening}
\textbf{By the end of this chapter:} \emph{I can talk about cleaning the house.}

Complete the last lesson of chapter 528: I can talk about cleaning the house.
\end{chapteropening}

\section[balayer]{balayer — to sweep}

\label{lesson:FR-C528-balayer}



\begin{quote}
Before the new one: say the French for wardrobe, then the French for to hang.
\end{quote}

\subsection*{You'll want to know: balayer}
\textbf{balayer} — \textquotedblleft{}to sweep\textquotedblright{}.

\begin{quote}
Tous les soirs, il balaie la cuisine.
\end{quote}

\begin{quote}
\emph{Every evening he sweeps the kitchen.}
\end{quote}

\emph{Je balaie} or \emph{je balaye}; both are correct.

\begin{cousinweb}
\emph{balayer} is built on \emph{balai} \textquotedblleft{}broom\textquotedblright{} (EtymDB), a word of Celtic origin, first the name of the broom plant.
\end{cousinweb}

\begin{grammarlens}[title={What you've built}]
One more word for telling, explaining and describing at B1.
\end{grammarlens}

\subsection*{Your turn}
\begin{itemize}
\raggedright
  \item \emph{Say it:} \emph{balayer}
  \item \emph{Say it:} \emph{balayer}, once more
  \item \emph{Say it:} a sentence of your own with \emph{balayer}
  \item \emph{Recall:} say the French for wardrobe, then the French for to hang, then say \emph{balayer} again
\end{itemize}

\begin{tcolorbox}[breakable,colback=teal!4,colframe=teal!35!black,title={Before you move on}]
What does balayer mean? (\textquotedblleft{}To sweep\textquotedblright{}.) Say it once more.
\end{tcolorbox}
\section[épousseter]{épousseter — to dust}

\label{lesson:FR-C528-epousseter}



\begin{quote}
Before the new one: say the French for to hang, then the French for to sweep.
\end{quote}

\subsection*{You'll want to know: épousseter}
\textbf{épousseter} — \textquotedblleft{}to dust\textquotedblright{}.

\begin{quote}
Elle époussette les meubles du salon.
\end{quote}

\begin{quote}
\emph{She dusts the furniture in the living room.}
\end{quote}

Spelling: \emph{j'époussette}. Built on \emph{la poussière}.

\begin{cousinweb}
\emph{épousseter} is built on an old word \emph{pousse} \textquotedblleft{}dust\textquotedblright{}, which goes back to Latin \emph{pulvis} \textquotedblleft{}dust\textquotedblright{} (Lewis \& Short); \emph{la poussière} has the same source.
\end{cousinweb}

\begin{grammarlens}[title={What you've built}]
One more word for telling, explaining and describing at B1.
\end{grammarlens}

\subsection*{Your turn}
\begin{itemize}
\raggedright
  \item \emph{Say it:} \emph{épousseter}
  \item \emph{Say it:} \emph{épousseter}, once more
  \item \emph{Say it:} a sentence of your own with \emph{épousseter}
  \item \emph{Recall:} say the French for to hang, then the French for to sweep, then say \emph{épousseter} again
\end{itemize}

\begin{tcolorbox}[breakable,colback=teal!4,colframe=teal!35!black,title={Before you move on}]
What does épousseter mean? (\textquotedblleft{}To dust\textquotedblright{}.) Say it once more.
\end{tcolorbox}
\section[frotter]{frotter — to rub, to scrub}

\label{lesson:FR-C528-frotter}



\begin{quote}
Before the new one: say the French for to sweep, then the French for to dust.
\end{quote}

\subsection*{You'll want to know: frotter}
\textbf{frotter} — \textquotedblleft{}to rub, to scrub\textquotedblright{}.

\begin{quote}
J'ai frotté la tache, mais elle est encore là.
\end{quote}

\begin{quote}
\emph{I scrubbed the stain, but it is still there.}
\end{quote}

\emph{Se frotter les mains} — to rub one's hands.

\begin{cousinweb}
\emph{frotter} goes back to Old French \emph{froter}. The origin is not settled; a link with Latin \emph{fricāre} \textquotedblleft{}to rub\textquotedblright{} is often proposed.
\end{cousinweb}

\begin{grammarlens}[title={What you've built}]
One more word for telling, explaining and describing at B1.
\end{grammarlens}

\subsection*{Your turn}
\begin{itemize}
\raggedright
  \item \emph{Say it:} \emph{frotter}
  \item \emph{Say it:} \emph{frotter}, once more
  \item \emph{Say it:} a sentence of your own with \emph{frotter}
  \item \emph{Recall:} say the French for to sweep, then the French for to dust, then say \emph{frotter} again
\end{itemize}

\begin{tcolorbox}[breakable,colback=teal!4,colframe=teal!35!black,title={Before you move on}]
What does frotter mean? (\textquotedblleft{}To rub, to scrub\textquotedblright{}.) Say it once more.
\end{tcolorbox}
\section[l'éponge]{l'éponge — sponge}

\label{lesson:FR-C528-eponge}



\begin{quote}
Before the new one: say the French for to dust, then the French for to rub.
\end{quote}

\subsection*{You'll want to know: l'éponge}
\textbf{l'éponge} — \textquotedblleft{}sponge\textquotedblright{}.

\begin{quote}
Il nettoie la table de la cuisine avec une éponge.
\end{quote}

\begin{quote}
\emph{He wipes down the kitchen table with a sponge.}
\end{quote}

Feminine. \emph{Passer l'éponge} — to let bygones be bygones.

\begin{cousinweb}
\emph{l'éponge} is inherited, through Old French \emph{esponge}, from Latin \emph{spongia} (EtymDB), a borrowing from Greek \emph{spongiá} \textquotedblleft{}a sponge\textquotedblright{}.
\end{cousinweb}

\begin{grammarlens}[title={What you've built}]
One more word for telling, explaining and describing at B1.
\end{grammarlens}

\subsection*{Your turn}
\begin{itemize}
\raggedright
  \item \emph{Say it:} \emph{l'éponge}
  \item \emph{Say it:} \emph{l'éponge}, once more
  \item \emph{Say it:} a sentence of your own with \emph{l'éponge}
  \item \emph{Recall:} say the French for to dust, then the French for to rub, then say \emph{l'éponge} again
\end{itemize}

\begin{tcolorbox}[breakable,colback=teal!4,colframe=teal!35!black,title={Before you move on}]
What does l'éponge mean? (\textquotedblleft{}Sponge\textquotedblright{}.) Say it once more.
\end{tcolorbox}
\section[la serpillière]{la serpillière — floor cloth, mop}

\label{lesson:FR-C528-serpilliere}



\begin{quote}
Before the new one: say the French for to rub, then the French for sponge.
\end{quote}

\subsection*{You'll want to know: la serpillière}
\textbf{la serpillière} — \textquotedblleft{}floor cloth, mop\textquotedblright{}.

\begin{quote}
Après le repas, je passe la serpillière.
\end{quote}

\begin{quote}
\emph{After the meal I mop the floor.}
\end{quote}

\emph{Passer la serpillière} — to mop. In the north of France and in Belgium also \emph{la wassingue}; in Québec \emph{la moppe}. Also spelt \emph{serpillère}.

\begin{cousinweb}
\emph{la serpillière} goes back to Old French \emph{sarpilliere} \textquotedblleft{}a coarse cloth\textquotedblright{}. The origin is not settled.
\end{cousinweb}

\begin{grammarlens}[title={What you've built}]
One more word for telling, explaining and describing at B1.
\end{grammarlens}

\subsection*{Your turn}
\begin{itemize}
\raggedright
  \item \emph{Say it:} \emph{la serpillière}
  \item \emph{Say it:} \emph{la serpillière}, once more
  \item \emph{Say it:} a sentence of your own with \emph{la serpillière}
  \item \emph{Recall:} say the French for to rub, then the French for sponge, then say \emph{la serpillière} again
\end{itemize}

\begin{tcolorbox}[breakable,colback=teal!4,colframe=teal!35!black,title={Before you move on}]
What does la serpillière mean? (\textquotedblleft{}Floor cloth, mop\textquotedblright{}.) Say it once more.
\end{tcolorbox}
